\chapter{Curvature Pinching and Preserved Curvature Properties}

\section{Curvature pinching and preserved curvature properties under Ricci flow}

\subsection{The Einstein tensor and its evolution}

\begin{definition}
\label{topping-ch9-einstein-tensor}
\source{topping:page-0106,topping:page-0107}
\dcref{ch9:2.1}
\group{curvature-pinching}
The Einstein tensor is the symmetric \((0,2)\)-tensor
\[
E=-\Ric+\frac{R}{2}g.
\]
In dimension three, with sectional curvatures
\(\lambda_1,\lambda_2,\lambda_3\) on an orthonormal basis of curvature
2-planes, the corresponding orthonormal basis diagonalises \(E\), whose
eigenvalues are \(\lambda_1,\lambda_2,\lambda_3\).
\end{definition}

\begin{proposition}
\label{topping-ch9-einstein-evolution}
\source{topping:page-0107}
\uses{topping-ch9-einstein-tensor}
\dcref{ch9:3.1,ch9:3.2}
\group{curvature-pinching}

Along Ricci flow,
\[
\left\langle \frac{\partial E}{\partial t}(X),W\right\rangle
=\Delta E(X,W)-2\langle \Rm(X,\cdot,W,\cdot),\Ric\rangle
+|\Ric|^2g(X,W).
\]
If the curvature operator is diagonal in an orthonormal frame, then
\[
\frac{\partial E}{\partial t}-\Delta E
=2\begin{pmatrix}
\lambda_1^2+\lambda_2\lambda_3 & 0 & 0\\
0 & \lambda_2^2+\lambda_1\lambda_3 & 0\\
0 & 0 & \lambda_3^2+\lambda_1\lambda_2
\end{pmatrix}.
\]
\end{proposition}

\subsection{The Uhlenbeck trick}

\begin{definition}
\label{topping-ch9-uhlenbeck-setup}
\source{topping:page-0108,topping:page-0109}
\dcref{ch9:4.1}
\group{curvature-pinching}
Let \((\mathcal M,g(t))\) be a Ricci flow and \(V\to\mathcal M\) a vector
bundle with an initial bundle isomorphism \(u_0:V\to T\mathcal M\) compatible
with the initial metric. Define \(u_t:V\to T\mathcal M\) by
\[
\frac{\partial u}{\partial t}=\Ric_{g(t)}(u),\qquad u(0)=u_0,
\]
and define the induced connection \(A(t)\) through
\[
u(A_Xv)=\nabla_X(u(v)).
\]
\end{definition}

\begin{remark}
\label{topping-ch9-uhlenbeck-metric}
\source{topping:page-0108,topping:page-0109}
\dcref{ch9:4.2,ch9:4.3}
\group{curvature-pinching}
Each \(u_t\) is a fibrewise isometry. Thus the induced fibre metric \(h\) on
\(V\) is time-independent, \(A(t)\) is compatible with \(h\), and
\(Au=Ah=0\).
\end{remark}

\begin{definition}
\label{topping-ch9-pulled-back-einstein}
\source{topping:page-0109,topping:page-0110}
\dcref{ch9:4.4,ch9:4.5}
\group{curvature-pinching}
The pull-back of the Einstein tensor is the section
\(\widetilde E\in\Gamma(V^*\otimes V)\) determined by
\[
u(\widetilde E(v))=E(u(v)).
\]
It is an \(h\)-symmetric endomorphism and obeys the evolution equation for
\(E\), with \(A(t)\) replacing \(\nabla\).
\end{definition}

\subsection{Parallel functions on vector bundles}

\begin{definition}
\label{topping-ch9-parallel-function}
\source{topping:page-0111}
\dcref{ch9:5.2}
\group{curvature-pinching}
Let \(W\to\mathcal M\) have fibre metric \(h\) and compatible connection
\(A(t)\). A smooth function \(\Psi:W\to\mathbb R\) is parallel when it is
invariant under parallel translation for every time.
\end{definition}

\begin{definition}
\label{topping-ch9-convex-parallel-function}
\source{topping:page-0112}
\dcref{ch9:5.3}
\group{curvature-pinching}
A parallel function \(\Psi:W\to\mathbb R\) is convex when its restriction
to each fibre \(W_p\) is convex.
\end{definition}

\begin{proposition}
\label{topping-ch9-parallel-composition-parabolic}
\source{topping:page-0112,topping:page-0113}
\uses{topping-ch9-parallel-function, topping-ch9-convex-parallel-function}
\dcref{ch9:5.4}
\group{curvature-pinching}

Suppose \(E(t)\in\Gamma(W)\) satisfies
\[
\frac{\partial E}{\partial t}-\Delta_AE=\Upsilon(E).
\]
For a parallel convex function \(\Psi\),
\[
\left(\frac{\partial}{\partial t}-\Delta_{\mathcal M}\right)(\Psi\circ E)
\le d\Psi(E)\!\left(\Upsilon(E)\right).
\]
\end{proposition}

\begin{definition}
\label{topping-ch9-parallel-convex-subset}
\source{topping:page-0113}
\dcref{ch9:5.5}
\group{curvature-pinching}
A subset \(X\subset W\) is parallel if parallel translation preserves it, and
it is convex if its intersection with every fibre is convex.
\end{definition}

\subsection{The ODE-PDE theorem}

\begin{theorem}
\label{topping-ch9-ode-pde-theorem}
\source{topping:page-0114,topping:page-0115,topping:page-0116}
\uses{topping-ch9-parallel-composition-parabolic, topping-ch9-parallel-convex-subset}
\dcref{ch9:6.1}
\group{curvature-pinching}

Let \(W\to(\mathcal M,g(t))\) have compatible connection \(A(t)\), and let
\(E(t)\in\Gamma(W)\) solve
\[
\frac{\partial E}{\partial t}-\Delta_AE=\Upsilon(E).
\]
If \(X\subset W\) is closed, convex, and parallel, and the fibrewise ODE
\[
\frac{de}{dt}=\Upsilon(e)
\]
preserves \(X\), then
\[
E(\mathord\cdot,0)\subset X\quad\Longrightarrow\quad
E(\mathord\cdot,t)\subset X\quad\text{for all }t\in[0,T].
\]
\end{theorem}

\begin{remark}
\label{topping-ch9-ode-pde-boundary-test}
\source{topping:page-0116}
\dcref{ch9:6.2}
\group{curvature-pinching}
If \(X\) is the zero sublevel set of a convex parallel function \(\Psi\), it
suffices to verify, for every ODE trajectory starting on \(\partial X\),
\[
\left.\frac{d}{dt}\Psi(e(t))\right|_{t=0}\le0.
\]
\end{remark}

\subsection{Applications to curvature pinching}

\begin{remark}
\label{topping-ch9-ode-eigenvalue-system}
\source{topping:page-0116,topping:page-0117}
\dcref{ch9:7.1}
\group{curvature-pinching}
For the pulled-back Einstein tensor, the fibrewise ODE leaves the eigenspaces
fixed and has eigenvalue system
\[
\frac{d\lambda_1}{dt}=2(\lambda_1^2+\lambda_2\lambda_3),\quad
\frac{d\lambda_2}{dt}=2(\lambda_2^2+\lambda_3\lambda_1),\quad
\frac{d\lambda_3}{dt}=2(\lambda_3^2+\lambda_1\lambda_2).
\]
\end{remark}

\begin{lemma}
\label{topping-ch9-eigenvalue-convexity}
\source{topping:page-0117,topping:page-0118}
\dcref{ch9:7.2}
\group{curvature-pinching}
For symmetric endomorphisms of \(\mathbb R^3\), the trace
\(\lambda_1+\lambda_2+\lambda_3\) is linear and the smallest eigenvalue
\(\lambda_1\) is concave. Hence \(\lambda_3\),
\(\lambda_1+\lambda_2\), and \(\lambda_3-\lambda_1\) are respectively
convex, concave, and convex. If \(\alpha\) is concave and \(\beta\) is
concave and increasing, then \(\beta\circ\alpha\) is concave wherever it is
defined.
\end{lemma}

\begin{theorem}
\label{topping-ch9-positive-sectional-preserved}
\source{topping:page-0118}
\uses{topping-ch9-eigenvalue-convexity, topping-ch9-ode-pde-theorem, topping-ch9-ode-eigenvalue-system}
\dcref{ch9:7.3}
\group{curvature-pinching}

For a closed three-dimensional manifold, weakly positive sectional curvature
is preserved by Ricci flow.
\end{theorem}

\begin{remark}
\label{topping-ch9-positive-sectional-strong}
\source{topping:page-0118}
\dcref{ch9:7.4}
\group{curvature-pinching}
The same maximum-principle argument, with a minor modification, preserves
strictly positive sectional curvature in the closed three-dimensional case.
\end{remark}

\begin{theorem}
\label{topping-ch9-ricci-pinching-preserved}
\source{topping:page-0119,topping:page-0120}
\uses{topping-ch9-eigenvalue-convexity, topping-ch9-ode-pde-theorem, topping-ch9-ode-eigenvalue-system}
\dcref{ch9:7.5}
\group{curvature-pinching}

For every \(\varepsilon\in[0,\tfrac13)\), the condition
\[
\Ric\ge\varepsilon Rg
\]
is preserved by Ricci flow on closed three-dimensional manifolds.
\end{theorem}

\begin{corollary}
\label{topping-ch9-ricci-positive-preserved}
\source{topping:page-0120}
\uses{topping-ch9-ricci-pinching-preserved}
\dcref{ch9:7.6,ch9:7.7}
\group{curvature-pinching}

On a closed three-dimensional manifold, Ricci flow preserves
\(\Ric>\varepsilon Rg\) for \(\varepsilon\in[0,\tfrac13)\), as well as
\(\Ric\ge0\) and \(\Ric>0\).
\end{corollary}

\begin{theorem}
\label{topping-ch9-roundness-estimate}
\source{topping:page-0121,topping:page-0122,topping:page-0123}
\uses{topping-ch9-eigenvalue-convexity, topping-ch9-ode-pde-theorem, topping-ch9-ode-eigenvalue-system, topping-ch9-ricci-pinching-preserved}
\dcref{ch9:7.8}
\group{curvature-pinching}

Given \(0<\beta<B<\infty\) and \(\gamma>0\), there is
\(M=M(\beta,B,\gamma)<\infty\) such that a Ricci flow on a closed
three-dimensional manifold satisfying
\[
\beta g(0)\le\Ric(g(0))\le Bg(0)
\]
also satisfies
\[
\left|\Ric-\frac13 Rg\right|\le\gamma R+M
\]
for every \(t\in[0,T]\).
\end{theorem}

\begin{remark}
\label{topping-ch9-roundness-reduction}
\source{topping:page-0121}
\dcref{ch9:7.9}
\group{curvature-pinching}
It is enough to control the nonnegative quantity \(\lambda_3-\lambda_1\),
because
\[
\left|\Ric-\frac13Rg\right|^2\le(\lambda_3-\lambda_1)^2.
\]
\end{remark}
